Sampling a discrete object from a prescribed, possibly non-uniform, distribution is a recurring primitive
in computational science, underlying Monte Carlo simulation, agent-based modeling, and the stochastic
components of modern machine learning pipelines, where billions of independent draws must be produced
under strict time and memory budgets. The related question of generating or sampling from a large,
implicitly described universe of discrete structures has likewise been studied extensively, from the
classical treatises of Nijenhuis and Wilf~\cite{NW75} and Kreher and Stinson~\cite{KS99} to Ruskey's
survey~\cite{Ruskey03}, more recent developments~\cite{GP21}, and Skiena's algorithmic
toolbox~\cite{Skiena20}. In practice, the objects to be generated are rarely equally likely: they carry
integer or rational \emph{weights}, and an unbiased sampler must reproduce the corresponding distribution
exactly, however skewed it is.

The reference solution to this weighted sampling problem is Walker's \emph{Alias
Method}~\cite{Walker74,Walker77}: it reorganizes a set of $n$ unequal probabilities into $n$ buckets of
equal capacity, each holding at most two outcomes, a primary value and its designated \emph{alias}, so that
a single bucket, drawn uniformly at random, together with one biased coin flip, suffices to recover the
target distribution. Vose later showed how to build this table in linear time~\cite{Vose91}, and Smith and
Jacobson provided a thorough analysis of the resulting sampler~\cite{SJ05}. By front-loading the whole
complexity of the distribution into a one-time linear-time preprocessing step, the Alias Method achieves
$O(1)$ sampling independently of $n$, making it the gold standard whenever a fixed distribution must be
queried a very large number of times. As for every sampler considered here, a one-time preprocessing is
followed by a sampling step repeated as often as needed, each with its own cost.

This efficiency, however, is traditionally obtained at the price of floating-point arithmetic, and modern
computational-science workloads increasingly expose the fragility of this choice. Both the construction of
the table and the sampling step rely on repeatedly comparing and renormalizing probabilities represented as
IEEE~754 doubles, so that rounding drift can silently corrupt the classification of outcomes during
preprocessing, and bit-creep can corrupt the coin-flip threshold during sampling. As soon as the total
weight of a distribution exceeds $2^{53}$, the granularity of double-precision floats becomes coarser than
the weight of individual items, which then acquire a \emph{structurally zero} probability of ever being
drawn. Such regimes are no longer isolated cases: in combinatorial chemistry, the space of chemically
plausible molecules grows so quickly with the number of atoms that it must be sampled rather than
enumerated, with reference libraries already ranging from about $10^{8}$ molecules for the smallest
structures to several trillion compounds for modern make-on-demand libraries~\cite{SSHL26}, well beyond
$2^{64}$ for molecules of about thirty heavy atoms. Monterrubio-Chanca and von Rudorff~\cite{MR25}, whose
open-source \texttt{nablachem} package~\cite{nablachem} we use in our own experiments, address this regime
with a pipeline whose final step is precisely a weighted random draw among combinatorial groups, over
weights of several hundred bits, well beyond the reach of a double-precision mantissa.

Several directions have been explored to make the Alias Method robust to this kind of scale. The most
direct remedy is to abandon floating-point numbers altogether and represent probabilities as
arbitrary-precision rational fractions (when it is possible), which restores mathematical exactness but incurs a substantial
overhead in both running time and memory, since every comparison and table entry now manipulates a pair of
big integers. An alternative, adopted by the Rust ecosystem's \texttt{weighted\_rand} crate~\cite{Rust25},
is to restrict the input to non-negative integer weights, instead of probabilities, so that the whole construction can be carried out
with exact integer arithmetic; this avoids fractions entirely, but its construction departs from Vose's and its correctness had, to the best
of our knowledge, never been established in the literature, and, as we show in \Cref{sec:algo}, it also produces
alias tables with systematically fewer alias-free cells than Vose's construction, which slows down the
sampling step. Orthogonally, entropy-optimal samplers such as \textsc{Fldr}~\cite{SFRM19} and its
refinement \textsc{Aldr}~\cite{DS26} build on the discrete-distribution-generating trees of Knuth and
Yao~\cite{KY76} to approach the information-theoretic minimum number of random bits per sample, at the cost
of a preprocessing and sampling cost that grows with the entropy of the distribution, again without native
support for arbitrary-precision weights.

In this paper we introduce the \emph{Integer-Alias Sampler}, a construction and sampling scheme that
operates natively and exclusively on arbitrary-precision integer weights, without ever resorting to
floating-point numbers or rational fractions. We revisit Walker and Vose's construction in this exact
integer setting, present the key-ideas for proving the correctness of the Rust integer-weight variant~\cite{Rust25}, and
prove that our own construction, built around a natural integer cell size, yields alias tables with more
alias-free cells, and hence a faster sampling step, together with matching tight bounds on the size of the
resulting table (\Cref{sec:walker,sec:algo}). We design an efficient random generation procedure on this
table, together with a virtual-object rejection mechanism, and bound its average random-bit consumption,
which we then compare experimentally with the near-optimal consumption of the entropy-optimal samplers
of~\cite{SFRM19,DS26}. We
implement all of the above under arbitrary-precision arithmetic, coupled with a background thread that
continuously refills a cryptographically secure random-bit stream, overlapping bit generation with sampling, supporting
distributions whose number of objects and individual weights can both grow without any fixed-width bound.
Finally, we experimentally compare our approach against the Fraction-Alias sampler, the Rust-Alias
sampler~\cite{Rust25}, and arbitrary-precision adaptations of \textsc{Fldr}~\cite{SFRM19} and
\textsc{Aldr}~\cite{DS26}, including a large-scale combinatorial-chemistry case study derived
from~\cite{MR25,SSHL26}, showing that once weights genuinely require arbitrary precision, our sampler offers
a favorable trade-off between sampling time, memory footprint, and random-bit consumption: it is faster than
the other exact Alias variants at comparable memory cost, while \textsc{Fldr} and \textsc{Aldr} consume fewer
random bits at a higher memory cost (\Cref{sec:exp}).
